\documentclass[11pt]{article}

\usepackage[T1]{fontenc}
\usepackage{lmodern}
\usepackage[a4paper,margin=27mm]{geometry}
\usepackage{microtype}
\usepackage{amsmath,amssymb,amsthm,mathtools,bm}
\usepackage{booktabs}
\usepackage[numbers,sort&compress]{natbib}
\usepackage[hidelinks]{hyperref}
\hypersetup{
  pdftitle={No history advantage in repeated environment-assisted qubit recovery},
  pdfauthor={Xiaojian Zeng}
}
\setlength{\emergencystretch}{2em}

\newtheorem{theorem}{Theorem}[section]
\newtheorem{lemma}[theorem]{Lemma}
\newtheorem{proposition}[theorem]{Proposition}
\newtheorem{corollary}[theorem]{Corollary}
\theoremstyle{remark}
\newtheorem{remark}[theorem]{Remark}

\newcommand{\cE}{\mathcal{E}}
\newcommand{\cM}{\mathcal{M}}
\newcommand{\cB}{\mathcal{B}}
\newcommand{\Tr}{\operatorname{tr}}
\newcommand{\rank}{\operatorname{rank}}
\newcommand{\id}{\operatorname{id}}
\newcommand{\R}{\mathbb{R}}
\newcommand{\C}{\mathbb{C}}
\newcommand{\one}{\mathbb{1}}
\newcommand{\dd}{\mathrm{d}}

\title{No history advantage in repeated environment-assisted qubit recovery}
\author{Xiaojian Zeng\\
\small PBC School of Finance, Tsinghua University, Beijing 100190, China\\
\small \texttt{xiaojianzeng@foxmail.com}}
\date{}

\begin{document}
\maketitle

\begin{abstract}
We compare current-outcome polar recovery with recovery conditioned on the
complete measurement history for repeated uses of a qubit channel.  The same
four-outcome projective measurement is made on each fresh minimal environment.
For every channel with full Choi rank and every positive use count, we prove
that the separately optimized entanglement fidelities coincide and equal
$(1+\lVert Q\rVert_*^n)/2$, where $Q$ polarizes the Kraus determinant.  The key
step is a quaternionic two-plane lemma that constructs, inside the complete
real Takagi-optimal family, a Kraus frame with pairwise commuting effects.
Thus retaining the outcome history gives no optimized fidelity advantage,
including at three uses where fixed noncommuting measurements can first have
complex branch traces.
\end{abstract}

\noindent\textbf{Keywords:} quantum feedback; environment-assisted error
correction; qubit channel; entanglement fidelity; Kraus decomposition; Takagi
factorization

\section{Introduction}

Measuring an environment after it interacts with a quantum system reveals a
classical error record.  Conditional operations on the system can then remove
the unitary part of the observed error.  This environment-assisted viewpoint
connects quantum feedback and error correction without introducing a code
subspace.  For a fixed fine-grained channel decomposition, polar recovery is
optimal at one step, and the optimal qubit recovery over several discrete
steps can be represented by one final polar correction
\cite{gregoratti2004}.  The latter representation makes the entire outcome
history available to the recovery.

Rexiti, Khabbazi Oskouei, and Mancini recently compared such history-dependent
(``Bayesian'') recovery with a memoryless (`Markovian') strategy that applies
the polar correction associated only with the current outcome
\cite{rexiti2024}.  In the version used here, equality is established for
Kraus-rank-two channels and for a special rank-three family, while selected
rank-three and rank-four cases are supported numerically.  Their Conjecture~9
asserts equality for every qubit channel and every number of uses.  The
remaining issue is not the optimal recovery for a fixed decomposition.  It is
whether an environment measurement can be chosen that is simultaneously
optimal for the history-dependent score and for memoryless polar recovery.

Three uses expose this issue sharply.  A trace of two positive $2\times2$
matrices is real, whereas a trace of three positive matrices can be complex.
Consequently, identities adequate at two uses do not by themselves close the
three-use comparison.  A generic environment measurement can indeed give a
strict pointwise history advantage.  A counterexample to the optimized
conjecture, however, would require a gap between two separately optimized
scores, not merely a gap at one nonoptimal measurement.

We prove that no such optimized gap exists for a qubit channel with
positive-definite Choi matrix.  More strongly, within the same fixed-measurement
and fresh-environment model, the optimized scores coincide for every positive
integer number of uses.  The proof has two ingredients.  First, Takagi
factorization identifies the global history-dependent optimum and leaves an
entire real orthogonal family of optimal Kraus frames.  Second, a new
four-dimensional quadratic-form lemma selects a member of this family whose
four effects commute.  The lemma uses quaternionic left--right multiplication
to find a two-plane on which three traceless quadratic forms are simultaneously
trace neutral.  Bases in that plane and its complement are then aligned by an
intersection argument.

The result is deliberately restricted.  The environment is fresh at each use,
the same orthogonal projective measurement is repeated, the outcome record is
perfect, and corrections are unitary polar recoveries.  Adaptive environment
measurements, additional-outcome positive-operator-valued measures, correlated
uses, imperfect records, higher system dimension, and nonunitary recovery
models are not covered.

\section{Channel, measurement, and feedback scores}
\label{sec:model}

Let $\cE:M_2(\C)\to M_2(\C)$ be completely positive and trace preserving
(CPTP), with positive-definite Choi matrix.  Its minimal Kraus rank is four.
Fix a minimal frame $A=(A_0,A_1,A_2,A_3)$,
\begin{equation}
  \cE(\rho)=\sum_{a=0}^{3}A_a\rho A_a^\dagger,
  \qquad
  \sum_{a=0}^{3}A_a^\dagger A_a=\one_2 .
  \label{eq:kraus}
\end{equation}
An orthogonal projective measurement of a minimal four-dimensional environment
is parametrized by $U\in U(4)$ and produces
\begin{equation}
  K_x(U)=\sum_{a=0}^{3}U_{xa}A_a,
  \qquad x=0,1,2,3 .
  \label{eq:mix}
\end{equation}
The same $U$ is used at every fresh, independent application of $\cE$.
Optimization over $U(4)$ removes dependence on the initially selected minimal
frame.

Write the polar decomposition as $K_x=V_xP_x$, where
$P_x=|K_x|=(K_x^\dagger K_x)^{1/2}\geq0$.  A partial isometry $V_x$ can be
extended unitarily on an unused kernel.  Applying $V_x^\dagger$ after the
current outcome gives the one-step corrected channel
\begin{equation}
  \cM_U(\rho)=\sum_xP_x\rho P_x,
  \qquad \sum_xP_x^2=\one_2 .
  \label{eq:markov-channel}
\end{equation}
After $n$ repetitions, its normalized entanglement fidelity is
\begin{equation}
 F^{\mathrm M}_n(U)
 =\frac14\sum_{x_1,\ldots,x_n}
 \left|\Tr(P_{x_n}\cdots P_{x_1})\right|^2 .
 \label{eq:markov-score}
\end{equation}
The modulus in Eq.~\eqref{eq:markov-score} is essential for $n\geq3$.

In the history-dependent protocol the correction at a given step may depend
on every outcome observed so far.  The optimal qubit score for a fixed
fine-grained decomposition is equivalently obtained by delaying the polar
recovery until the end \cite{gregoratti2004}.  In the present notation it is
\begin{equation}
 F^{\mathrm B}_n(U)
 =\frac14\sum_{x_1,\ldots,x_n}
 \left(\Tr\left|K_{x_n}\cdots K_{x_1}\right|\right)^2 .
 \label{eq:bayes-score}
\end{equation}
We compare the separately optimized quantities
\begin{equation}
 F^{\mathrm M,*}_n(\cE)=\max_{U\in U(4)}F^{\mathrm M}_n(U),
 \qquad
 F^{\mathrm B,*}_n(\cE)=\max_{U\in U(4)}F^{\mathrm B}_n(U).
 \label{eq:optimized}
\end{equation}
Because current-outcome polar recovery is an admissible history-dependent
policy, $F^{\mathrm M}_n(U)\leq F^{\mathrm B}_n(U)$ for each fixed $U$.

\section{Determinant reduction and the three-use obstruction}
\label{sec:reduction}

For every complex $2\times2$ matrix $X$,
\begin{equation}
  (\Tr|X|)^2=\Tr(X^\dagger X)+2|\det X| .
  \label{eq:two-singular-values}
\end{equation}
Summing the first term over all length-$n$ Kraus histories gives
$\Tr\one_2=2$, while determinants factor along a history.  Hence, with
\begin{equation}
  s(U)=\sum_x|\det K_x(U)|,
  \label{eq:s-def}
\end{equation}
the fixed-measurement Bayesian score reduces exactly to
\begin{equation}
  F^{\mathrm B}_n(U)=\frac{1+s(U)^n}{2}.
  \label{eq:bayes-reduction}
\end{equation}

The determinant of a linear combination of $2\times2$ matrices is a
homogeneous quadratic form.  There is therefore a complex symmetric matrix
$Q=Q^{\mathsf T}\in M_4(\C)$ such that
\begin{equation}
  \det\!\left(\sum_a z_aA_a\right)=z^{\mathsf T}Qz .
  \label{eq:q-def}
\end{equation}
Equation~\eqref{eq:mix} gives
\begin{equation}
  s(U)=\sum_x |(UQU^{\mathsf T})_{xx}|
  \leq \|UQU^{\mathsf T}\|_*=\|Q\|_* ,
  \label{eq:pinching}
\end{equation}
where $\|Q\|_*=\Tr\sqrt{Q^\dagger Q}$ is the trace norm.  If
$Q=V\Sigma V^{\mathsf T}$ is a Takagi factorization with
$\Sigma=\operatorname{diag}(\sigma_0,\ldots,\sigma_3)\geq0$, then
$U=V^\dagger$ attains the bound.  Thus
\begin{equation}
 F^{\mathrm B,*}_n(\cE)=\frac{1+\|Q\|_*^n}{2}.
 \label{eq:bayes-optimum}
\end{equation}
Takagi decompositions play an analogous determinant-ensemble role in
two-qubit concurrence problems \cite{laustsen2003}; Eq.~\eqref{eq:bayes-optimum}
is not the missing part of the feedback comparison.

To see what is missing at three uses, expand each positive effect square root
as
\begin{equation}
  P_x=a_x\one_2+\bm b_x\cdot\bm\sigma,
  \qquad a_x\in\R,\quad \bm b_x\in\R^3,
\end{equation}
and set $G=\sum_x\bm b_x\bm b_x^{\mathsf T}\geq0$.  Trace preservation in
Eq.~\eqref{eq:markov-channel} implies
\begin{equation}
  s(U)=1-2\Tr G,
  \qquad R(U)=s(U)I_3+2G,
  \label{eq:bloch-map}
\end{equation}
where $R(U)$ is the Bloch matrix of $\cM_U$.  Consequently
\begin{equation}
  F^{\mathrm M}_n(U)=\frac{1+\Tr R(U)^n}{4}.
  \label{eq:markov-bloch}
\end{equation}
Let $\beta_1,\beta_2,\beta_3\geq0$ be the eigenvalues of $G$,
$t=\sum_i\beta_i=(1-s)/2$,
$e_2=\sum_{i<j}\beta_i\beta_j$, and
$e_3=\beta_1\beta_2\beta_3$.  Direct expansion yields
\begin{equation}
 F^{\mathrm B}_3(U)-F^{\mathrm M}_3(U)
 =6\big[(s+t)e_2-e_3\big]\geq0.
 \label{eq:three-deficit}
\end{equation}
Apart from the noiseless boundary, equality holds precisely when
$\rank G\leq1$, or equivalently when all $P_x$ commute.  Thus a three-use
proof requires a Kraus frame that both attains Eq.~\eqref{eq:pinching} and has
commuting effects.  Takagi factorization alone does not provide the second
property.

\section{A collinear-diagonal basis for three quadratic forms}
\label{sec:lemma}

We now prove the geometric statement that supplies the missing compatibility.
It is distinct from simultaneous hollowization of a prescribed pair of
matrices \cite{damm2020}: the output line is free to be chosen, and three
traceless forms are treated in the special dimension four.

\begin{lemma}[Collinear diagonals]
\label{lem:collinear}
Let $H_1,H_2,H_3$ be real symmetric $4\times4$ matrices with zero trace.
There is an orthonormal basis $w_1,\ldots,w_4$ of $\R^4$ such that the four
vectors
\begin{equation}
 q(w_j)=\bigl(w_j^{\mathsf T}H_1w_j,
                    w_j^{\mathsf T}H_2w_j,
                    w_j^{\mathsf T}H_3w_j\bigr)\in\R^3
 \label{eq:q-vector}
\end{equation}
are collinear.  Zero vectors are allowed.
\end{lemma}

\begin{proof}
Identify $\R^4$ with the quaternions $\mathbb H$ in the orthonormal basis
$(1,i,j,k)$.  For pure unit quaternions $u,v$, denote left and right
multiplication by $L_u$ and $R_v$.  These real matrices are skew symmetric,
square to $-I_4$, and commute with one another.  Hence
\begin{equation}
  \Gamma(u,v)=L_uR_v
\end{equation}
is a symmetric orthogonal involution.  It has zero trace, so
\begin{equation}
  P(u,v)=\frac{I_4+\Gamma(u,v)}{2}
  \label{eq:projection}
\end{equation}
is a rank-two orthogonal projection.

Put $(e_1,e_2,e_3)=(i,j,k)$.  For a traceless symmetric $H$, define the real
$3\times3$ matrix
\begin{equation}
 C(H)_{ab}=\Tr\bigl(HL_{e_a}R_{e_b}\bigr).
 \label{eq:c-matrix}
\end{equation}
Writing pure quaternions in their three real coordinates and using
$\Tr H=0$ gives
\begin{equation}
 \Tr\bigl(P(u,v)H\bigr)=\frac12u^{\mathsf T}C(H)v .
 \label{eq:neutral-identity}
\end{equation}
For $C_k=C(H_k)$, consider the continuous cubic
\begin{equation}
 g(u)=\det\big[C_1^{\mathsf T}u,C_2^{\mathsf T}u,C_3^{\mathsf T}u\big],
 \qquad u\in S^2 .
 \label{eq:odd-cubic}
\end{equation}
It is odd.  A path from $u$ to $-u$ therefore contains a zero $u_*$.  At
that zero the three vectors $C_k^{\mathsf T}u_*$ span a subspace of dimension
at most two.  Choose a unit $v_*$ perpendicular to their span.  By
Eq.~\eqref{eq:neutral-identity},
\begin{equation}
 \Tr(PH_k)=0,\qquad k=1,2,3,
 \label{eq:neutral-plane}
\end{equation}
for $P=P(u_*,v_*)$.  Since $\Tr H_k=0$, the same identities hold on
$P^\perp$.  We have therefore split $\R^4$ into two orthogonal two-planes on
which all three compressed forms are trace neutral.

Choose an orthonormal basis $(p_1,p_2)$ in the range of $P$.  Each compressed
form has the shape
\begin{equation}
 \begin{pmatrix}a_k&b_k\\ b_k&-a_k\end{pmatrix}.
\end{equation}
For $p(\theta)=\cos\theta\,p_1+\sin\theta\,p_2$,
\begin{equation}
 q(p(\theta))=\bm a\cos(2\theta)+\bm b\sin(2\theta),
 \label{eq:ellipse}
\end{equation}
where $\bm a=(a_1,a_2,a_3)$ and
$\bm b=(b_1,b_2,b_3)$.  Its orthogonal partner has the negative $q$ vector.
The same description holds on $P^\perp$.  Each pair therefore traces a
centered ellipse contained in a linear subspace of $\R^3$ of dimension at
most two.

If either ellipse map has rank at most one, its $3\times2$ defining map has a
nonzero kernel.  Choosing the corresponding angle makes both $q$ vectors in
that two-plane zero; any orthonormal pair in the other plane gives an opposite
pair, so all four vectors are collinear.  If both maps have rank two, their
two-dimensional ranges in $\R^3$ intersect in a line.  Choose angles whose
images lie on that common line.  The orthogonal partners again give the
negatives.  These cases include every possible degeneracy and prove the
claim.
\end{proof}

\section{A Takagi-optimal frame with commuting effects}
\label{sec:construction}

The real freedom left after Takagi factorization is exactly large enough to
apply Lemma~\ref{lem:collinear}.  Start from
$Q=V\Sigma V^{\mathsf T}$ and mix the original frame by $V^\dagger$, writing
the resulting Kraus operators as $B_0,\ldots,B_3$.  For a real vector
$r\in\R^4$, put
\begin{equation}
  K(r)=\sum_{j=0}^{3}r_jB_j .
\end{equation}
Then
\begin{equation}
  \det K(r)=r^{\mathsf T}\Sigma r\geq0.
  \label{eq:positive-det}
\end{equation}
Consequently, for every $O\in O(4)$, the four rows $r_x$ define a frame mixed
by $U=OV^\dagger$ for which
\begin{equation}
 \sum_x|\det K(r_x)|
 =\sum_{x,j}\sigma_jO_{xj}^2
 =\Tr\Sigma=\|Q\|_* .
 \label{eq:real-takagi-family}
\end{equation}
Every member of this real orthogonal family is globally optimal for the
Bayesian score.

Expand the effect of $K(r)$ in the Pauli basis:
\begin{equation}
 K(r)^\dagger K(r)
 =\alpha(r)\one_2+\sum_{k=1}^{3}(r^{\mathsf T}H_kr)\sigma_k .
 \label{eq:effect-forms}
\end{equation}
Because $r$ is real, each coefficient is a real quadratic form and the
$H_k$ may be chosen real symmetric.  More explicitly,
\begin{equation}
 r^{\mathsf T}H_kr=\frac12\Tr\bigl[\sigma_kK(r)^\dagger K(r)\bigr].
\end{equation}
Completeness of the Takagi-mixed frame gives
\begin{equation}
 \Tr_{\R^4}H_k
 =\frac12\Tr\!\left[\sigma_k\sum_jB_j^\dagger B_j\right]
 =\frac12\Tr\sigma_k=0 .
 \label{eq:traceless-forms}
\end{equation}
Apply Lemma~\ref{lem:collinear} to $(H_1,H_2,H_3)$ and use its basis vectors
as the rows of $O$.  Equation~\eqref{eq:effect-forms} then says that the Bloch
vectors of the four effects $K(r_x)^\dagger K(r_x)$ lie on one line.  Two
qubit Hermitian matrices commute precisely when their Bloch vectors are
collinear.  The four effects, and hence their positive square roots, commute.
Combined with Eq.~\eqref{eq:real-takagi-family}, this proves the following.

\begin{proposition}[Optimal commuting frame]
\label{prop:commuting-frame}
Every full-Choi-rank qubit CPTP map has a four-outcome orthogonal environment
measurement that simultaneously satisfies
\begin{enumerate}
 \item $\sum_x|\det K_x|=\|Q\|_*$, and
 \item $[K_x^\dagger K_x,K_y^\dagger K_y]=0$ for all $x,y$.
\end{enumerate}
\end{proposition}

No nondegeneracy of the Takagi singular values or of the individual outcome
operators was used.  In particular, zero or singular outcomes do not obstruct
the construction.

\section{Equality for every use count}
\label{sec:main}

\begin{theorem}[No optimized history advantage]
\label{thm:main}
Let $\cE$ be a qubit CPTP map with positive-definite Choi matrix.  Use the
same four-outcome orthogonal projective measurement on every fresh minimal
environment.  For every positive integer $n$, the separately optimized
entanglement fidelities of current-outcome polar recovery and full-history
optimal recovery satisfy
\begin{equation}
 \boxed{
 F^{\mathrm M,*}_n(\cE)=F^{\mathrm B,*}_n(\cE)
 =\frac{1+\|Q\|_*^n}{2}.}
 \label{eq:main-result}
\end{equation}
In particular, the optimized three-use gap is zero.
\end{theorem}

\begin{proof}
Choose the frame from Proposition~\ref{prop:commuting-frame} and diagonalize
its commuting square roots in one system basis:
\begin{equation}
 P_x=|K_x|=\begin{pmatrix}p_x&0\\0&q_x\end{pmatrix},
 \qquad p_x,q_x\geq0 .
\end{equation}
Trace preservation and determinant optimality give
\begin{equation}
 \sum_xp_x^2=\sum_xq_x^2=1,
 \qquad \sum_xp_xq_x=\sum_x|\det K_x|=\|Q\|_* .
 \label{eq:diagonal-sums}
\end{equation}
The corrected channel $\cM_U$ preserves the longitudinal Bloch component and
multiplies both transverse components by $\|Q\|_*$.  Its Bloch eigenvalues are
therefore $(1,\|Q\|_*,\|Q\|_*)$, and Eq.~\eqref{eq:markov-bloch} gives
\begin{equation}
 F^{\mathrm M}_n(U)=\frac{1+\|Q\|_*^n}{2}.
\end{equation}
By Eqs.~\eqref{eq:bayes-reduction} and
\eqref{eq:real-takagi-family}, the Bayesian score at the same $U$ has exactly
this value, which is its global optimum in Eq.~\eqref{eq:bayes-optimum}.
Finally, $F^{\mathrm M,*}_n\leq F^{\mathrm B,*}_n$ because the memoryless
policy is a special history-dependent policy, whereas the constructed frame
shows the reverse inequality at the common optimum.  This proves
Eq.~\eqref{eq:main-result}.
\end{proof}

The theorem does not say that history is useless at a prescribed environment
measurement.  Equation~\eqref{eq:three-deficit} shows that a noncommuting
fixed frame can have a strict three-use advantage.  Rather, it says that an
optimal measurement can always be chosen from the Takagi family so that the
memoryless corrected channel already reaches the global history-dependent
ceiling.  This separation between pointwise and optimized statements is
essential to the physical interpretation.

\section{Deterministic numerical checks}
\label{sec:checks}

The argument above is finite dimensional and exact; numerical evidence is not
used to infer Theorem~\ref{thm:main}.  Deterministic scripts accompanying this
article implement the quaternionic construction and independently enumerate
Kraus histories as diagnostic checks.  Twelve arbitrary triples of traceless
real symmetric forms and eight nonunital channels with strictly positive Choi
spectra were tested.  The latter used $n=1,\ldots,7$ transfer calculations and
an explicit $4^3=64$-history calculation at $n=3$.  Table~\ref{tab:checks}
reports the largest residual in each class.

\begin{table}[ht]
\centering
\caption{Floating-point diagnostics for the exact construction.  These
residuals check the implementation and are not substitutes for the proof.}
\label{tab:checks}
\begin{tabular}{lr}
\toprule
Diagnostic & largest residual \\
\midrule
generic three-form construction & $1.639\times10^{-15}$ \\
Takagi diagonalization & $1.364\times10^{-15}$ \\
determinant optimum $s=\|Q\|_*$ & $4.441\times10^{-16}$ \\
effect commutators & $1.316\times10^{-16}$ \\
score identity, $n=1,\ldots,7$ & $1.554\times10^{-15}$ \\
explicit 64-history identity, $n=3$ & $3.331\times10^{-15}$ \\
\bottomrule
\end{tabular}
\end{table}

The smallest Choi eigenvalue among these eight channel instances was
$3.662\times10^{-3}$.  The same package reproduces the published special
rank-three baseline of Ref.~\cite{rexiti2024} for three interior parameters
and $n=1,2,3,5$, with largest residual $1.887\times10^{-15}$.  It also checks
the one-step polar-recovery result of Ref.~\cite{gregoratti2004} and verifies
the complex-modulus convention in Eq.~\eqref{eq:markov-score} by a direct
normalized Choi--Bell contraction.

\section{Discussion and limitations}
\label{sec:discussion}

The trace norm of the determinant polarization $Q$ controls the common
optimum, but the main structural result is not the Takagi bound itself.  The
nontrivial step is compatibility: among all determinant-optimal measurements,
there is one whose four effects commute.  Lemma~\ref{lem:collinear} establishes
this without assuming a unital channel, distinct Takagi singular values, or
invertible outcome operators.  It also explains why a generic simultaneous
hollowization result for a fixed pair of matrices is insufficient: the proof
must exploit both four real Kraus coordinates and the freedom to select the
surviving common Bloch axis.

At three uses, noncommuting positive factors may have a complex trace, and a
fixed measurement can make full outcome history operationally useful.  The
optimization over the projective environment basis nevertheless removes this
advantage.  The all-$n$ extension follows because the constructed channel has
the exact Bloch spectrum $(1,s,s)$; it is not a perturbative extrapolation from
the three-use case.

Several nearby questions remain open outside the theorem's domain.  An
adaptive choice of environment basis changes the admissible optimization.
Extra-outcome nonprojective measurements are not parametrized by $U(4)$.
Correlated environments invalidate the product-history factorization behind
Eq.~\eqref{eq:bayes-reduction}.  Nonunitary correcting channels, imperfect
records, and higher-dimensional systems require different optimality
arguments.  No experimental performance or implementation-cost claim is made.
Finally, the literature comparison supporting priority is bounded rather than
exhaustive and requires human review before submission.

\section*{Statements and declarations}

\paragraph{Data and code availability.}
There are no experimental data.  The reproducibility package accompanying
this draft contains the Python source, fixed-seed logs, environment
requirements, and commands for the published-baseline and construction
checks described in Sec.~\ref{sec:checks}.

\paragraph{Funding.}
No funding information was supplied during preparation of this draft.  The
author must confirm or replace this statement before submission.

\paragraph{Competing interests.}
No competing-interest declaration was supplied during preparation of this
draft.  The author must provide and approve the final declaration before
submission.

\paragraph{Use of generative artificial intelligence.}
OpenAI Codex was used on 12--13 September 2026 to propose and develop the
analytic argument, derive and check equations, write verification code,
conduct a bounded literature comparison, perform AI-only scientific checking,
and draft the manuscript.  AI-only checking is not human peer review.  The
listed author has not yet completed independent verification of the argument,
code, sources, or final text; this draft must not be submitted until the author
has done so and accepts scholarly responsibility for the work.

\bibliographystyle{unsrtnat}
\bibliography{references}

\end{document}
